\documentclass{article}

% -------------------------------------------------
% Packages
% -------------------------------------------------

\usepackage{xcolor}
\usepackage{amsmath}
\usepackage{amssymb}
\usepackage{amsthm}
\usepackage{graphicx}
\usepackage{subcaption}
\usepackage{matlab-prettifier}
\usepackage{listings}
\usepackage{blindtext}
\usepackage{geometry}
\usepackage{float}
\usepackage{dirtytalk}
\usepackage{enumitem}
\usepackage{derivative}
\usepackage{biblatex}
\usepackage{svg}
\usepackage{hyperref}

% -------------------------------------------------
% Colors
% -------------------------------------------------

\definecolor{darkgreen}{rgb}{0.0, 0.5, 0.0}

% -------------------------------------------------
% Hyperlinks
% -------------------------------------------------

\hypersetup{
    colorlinks=true,
    linkcolor=blue,
    citecolor=blue,
    urlcolor=blue
}

% -------------------------------------------------
% Theorem environments
% -------------------------------------------------

% Numbered
\newtheorem{theorem}{Theorem}[section]
\newtheorem{lemma}[theorem]{Lemma}
\newtheorem{proposition}[theorem]{Proposition}

% Unnumbered
\newtheorem*{theorem*}{Theorem}
\newtheorem*{lemma*}{Lemma}
\newtheorem*{proposition*}{Proposition}

\theoremstyle{definition}
\newtheorem{definition}[theorem]{Definition}
\newtheorem{example}[theorem]{Example}

\theoremstyle{remark}
\newtheorem{remark}[theorem]{Remark}

% -------------------------------------------------
% Equation numbering
% -------------------------------------------------

% Reset equation numbering at each subsection
\makeatletter
\@addtoreset{equation}{subsection}
\makeatother

% If you want equations to be numbered by subsection,
% uncomment the following line:
% \renewcommand{\theequation}{\thesubsection.\arabic{equation}}

% -------------------------------------------------
% Python / MATLAB code formatting
% -------------------------------------------------

\lstset{
    language=Python,
    basicstyle=\ttfamily\small,
    keywordstyle=\color{blue},
    stringstyle=\color{darkgreen},
    commentstyle=\color{gray},
    showstringspaces=false,
    frame=single,
    numbers=left,
    numberstyle=\tiny\color{gray},
    breaklines=true
}

% -------------------------------------------------
% Page geometry
% -------------------------------------------------

\geometry{
    a4paper,
    total={160mm,250mm},
    left=25mm,
    top=30mm,
    footskip=10mm
}

% -------------------------------------------------
% Document settings
% -------------------------------------------------

\setlength{\parindent}{0pt}

% -------------------------------------------------
% Title
% -------------------------------------------------

\title{Calculus Uge 2}
\author{Thomas Sejer Canham}
\date{\today}

% =================================================
% Document
% =================================================

\begin{document}

% -------------------------------------------------
% Title page
% -------------------------------------------------

\begin{titlepage}

    \maketitle


I denne uge skal vi arbejde med integration ved substitution, delvis integration og vi vil kigge på hvordan i sammenspil med den fundamentale sætning for calculus kan bruges til at beregne arealer under kurver. Bemærk, at integraler er generelt svære at beregne, så dette er nogle af de mest anvendte analytiske metoder, der kan hjælpe os med at finde løsninger.
\section*{Integration ved substitution}
Integration ved substitution vil typisk blive brugt, når vi har en funktion, der er sammensat af to funktioner. For eksempel, hvis vi har en funktion $f(g(x))$, hvor $g(x)$ er en indre funktion og $f$ er en ydre funktion, kan vi bruge substitution til at forenkle integralet. Den bagvedliggende ide er at koordinattrasformationen $u = g(x)$, frembringer en gunstig form på integralet, hvor $\dfrac{dg}{dx}$ (Jacobian) kan bruges til at erstatte $dx$ i integralet. Dette kan gøre integralet lettere at evaluere.
\subsection*{Eksempel}
Vi er givet integralet:
\begin{equation}
    \int 2x \cos(x^2) dx
\end{equation}

Vi bruger substitution: $g(x) = u = x^2$, derved bliver $\dfrac{du}{dx} = 2x$, så $du = 2x \, dx$. Kigger vi nærmere på ovenstående integrale, kan vi se at netop $2x \, dx$ er til stede, hvilket gør substitutionen mulig. Vi kan nu omskrive integralet som:
\begin{equation}
    \int \cos(u) du = \sin(u) + C = \sin(x^2) + C
\end{equation}

\section*{Fundamentale sætning for calculus}
Den fundamentale sætning for calculus siger, at hvis $f$ er en kontinuert funktion på intervallet $[a, b]$, og vi beregner
\begin{equation}
    F(x) = \int_a^x f(t) dt
\end{equation}
så er $F'(x) = f(x)$ for alle $x$ i intervallet $[a, b]$. Dette definerer en stamfunktion $F$ til $f$. Ydermere, hvis $F$ er en stamfunktion til $f$, så gælder det, at
\begin{equation}
    \int_a^b f(x) dx = F(b) - F(a)
\end{equation}
Dette giver os en metode til at beregne bestemte integraler ved at finde stamfunktioner. 


\end{titlepage}



\section*{Delvis integration}
Delvis integration en en metode der er motiveret af produktreglen for differentiation. Hvis vi har to funktioner $f(x)$ og $g(x)$, og lader $h(x) = f(x)g(x)$, så gælder det at
\begin{equation}
    \frac{d}{dx} h(x) = h'(x) = f'(x)g(x) + f(x)g'(x)
\end{equation}
Bruger vi nu fundamentale sætning for calculus, kan vi integrere begge sider af ligningen:
\begin{equation}
    h(x) = \int h'(x)  dx = [g(x)f(x)] = \int f'(x)g(x) dx + \int f(x)g'(x) dx
\end{equation}
der af følger det at
\begin{equation}
    \int f(x)g'(x) dx = [f(x)g(x)] - \int f'(x)g(x) dx.
    \end{equation}  
Måden dette skal forstås er, at vi har et udtryk, hvor vi kan vælge at definere $g'(x)$ og $f(x)$ på en sådan måde, at vi forhåbenligt kender stamfunktionen $g(x)$, således at $[f(x)g(x)]$ er kendt og at integralet $\int f'(x)g(x) dx$ er lettere at beregne end det oprindelige integral $\int f(x)g'(x) dx$.
\subsection*{Eksempel}
En klassisk anvendelse af delvis integration er at beregne integralet
\begin{equation}
    \int x e^x dx
\end{equation}
Vi kan vælge at lade $f(x) = x$ og $g'(x) = e^x$. Vi kender stamfunktionen til $g'(x)$, som er $g(x) = e^x$, og vi kan differentiere $f(x)$ for at få $f'(x) = 1$. Anvender vi delvis integration, får vi:
\begin{equation}
    \int x e^x dx = [x e^x] - \int 1 \cdot e^x dx = x e^x - e^x + C = e^x(x - 1) + C
    \end{equation}



\end{document}